\documentclass[a4paper, serif, 10pt]{moderncv}

%% moderncv
% color and style
\moderncvcolor{black}
\moderncvstyle{banking}

%% page layout/margins
\usepackage[top=2.5cm, bottom=2.5cm, left=1.5cm, right=1.5cm]{geometry}

\nopagenumbers{}

%% Babel config for Brazilian Portuguese language
% ref:
% - https://latex3.github.io/babel/guides/locale-portuguese.html
\usepackage[brazilian]{babel}

%% fontspec
\usepackage{fontspec}

\IfFontExistsTF{Linux Libertine}{
  \setmainfont[Ligatures={TeX, Common}, Numbers=OldStyle]{Linux Libertine}
}{}
\IfFontExistsTF{Linux Libertine O}{
  \setmainfont[Ligatures={TeX, Common}, Numbers=OldStyle]{Linux Libertine O}
}{}

%% CTeX
% CJK support, used to show CJK characters in the resume
%
% - heading=false: disable ctex heading
% - punct=kaiming: use kaiming punctuations styles for CJK
% - scheme=plain: use plain scheme, do not override \normalsize font size
% - space=auto: space settings for CJK characters
%
% CJK fonts are resolved per locale language, e.g. Noto Serif/Sans CJK JP
% for Japanese resumes, with fallback to the next region if not installed
%
% ref:
% - http://ctan.mirrorcatalogs.com/language/chinese/ctex/ctex.pdf
\usepackage[UTF8, heading=false, punct=kaiming, scheme=plain, space=auto]{ctex}

\IfFontExistsTF{Noto Serif CJK SC}{
  \setCJKmainfont{Noto Serif CJK SC}
}{}
\IfFontExistsTF{Noto Sans CJK SC}{
  \setCJKsansfont{Noto Sans CJK SC}
}{}

%% line spacing
\usepackage{setspace}
\setstretch{1.125}

%% URL styling - use normal text instead of monospace
\urlstyle{same}

%% Patch \httplink and \httpslink to handle full URLs
% The original moderncv macros always prepend http:// or https:// to URLs.
% This causes issues when the URL already has a protocol (e.g., https://example.com)
% resulting in malformed URLs like https://https://example.com.
% This patch checks if the URL already contains "://" and uses it directly if so.
% Note: We use \str_set:Nx (x-expansion) to expand macros like \@homepage first.
\ExplSyntaxOn
\str_new:N \g_yamlresume_url_str

\RenewDocumentCommand{\httpslink}{O{}m}{
  \str_set:Nx \g_yamlresume_url_str {#2}
  \str_if_in:NnTF \g_yamlresume_url_str {://}
    { \url{#2} }
    { \url{https://#2} }
}

\RenewDocumentCommand{\httplink}{O{}m}{
  \str_set:Nx \g_yamlresume_url_str {#2}
  \str_if_in:NnTF \g_yamlresume_url_str {://}
    { \url{#2} }
    { \url{http://#2} }
}
\ExplSyntaxOff

\name{Mariana Costa Albuquerque}{}
\title{Engenheira Mecânica Sênior | Projetos, Simulação e Manufatura}
\phone[mobile]{+55 11 98765-4321}
\email{mariana.albuquerque@engmail.com.br}
\homepage{https://www.marianaalbuquerque.eng.br}

\address{Av. Paulista, 1578, Bela Vista, São Paulo, São Paulo, Brasil, 01310-100}{}{}

\extrainfo{{\small \faLinkedin}\ \href{https://www.linkedin.com/in/mariana-albuquerque}{@mariana-albuquerque} {} {} {} • {} {} {} 
{\small \faGithub}\ \href{https://github.com/mariana-albuquerque}{@mariana-albuquerque} {} {} {} • {} {} {} 
{\small \faStackOverflow}\ \href{https://stackoverflow.com/users/12345678/mariana-albuquerque}{@mariana-albuquerque}}

\begin{document}

\maketitle

\section{Informações Básicas}

\cvline{}{Engenheira mecânica com mais de 10 anos de experiência em projeto, simulação e industrialização de sistemas mecânicos para os setores automotivo, energético e de bens de capital.
%
\begin{itemize}
\item Especialista em CAD 3D, análise estrutural por elementos finitos e otimização de componentes.
\item Histórico comprovado na redução de custos, melhoria de desempenho e transferência de produtos para a manufatura.
\item Experiência em gestão de projetos multidisciplinares, normas ISO 9001 e metodologias Lean Six Sigma.
\item Perfil colaborativo, com forte comunicação técnica e mentoria de equipes de engenharia.
\end{itemize}}
\section{Formação}

\cventry{mar. de 2012–fev. de 2014}
        {Mestrado, Engenharia Mecânica, Nota: 3.9/4.0}
        {Universidade Estadual de Campinas}
        {\url{https://www.unicamp.br}}
        {}
        {\begin{itemize}
\item Dissertação voltada à otimização topológica de suportes estruturais, com validação experimental e simulação numérica.
\item Desenvolvimento de métodos para reduzir massa mantendo rigidez e confiabilidade em aplicações automotivas.
\item Publicação de resultados em congresso nacional de engenharia mecânica.
\end{itemize}
\textbf{Cursos}: Mecânica dos Sólidos Computacional, Dinâmica de Sistemas Mecânicos, Tribologia e Lubrificação, Materiais Compósitos, Metodologia de Pesquisa Experimental, Otimização Estrutural}

\cventry{mar. de 2006–dez. de 2010}
        {Bacharelado, Engenharia Mecânica, Nota: 3.7/4.0}
        {Universidade Federal de Santa Catarina}
        {\url{https://www.ufsc.br}}
        {}
        {\begin{itemize}
\item Formação sólida em projeto mecânico, materiais e processos de fabricação.
\item Participação em equipe de competição universitária com foco em eficiência energética.
\item Experiência prática em laboratórios de ensaios mecânicos e metrologia.
\end{itemize}
\textbf{Cursos}: Cálculo Diferencial e Integral, Resistência dos Materiais, Termodinâmica Aplicada, Mecânica dos Fluidos, Processos de Fabricação, Desenho Técnico Mecânico, Controle de Sistemas Dinâmicos}
\section{Experiência Profissional}

\cventry{jan. de 2021–Atual}
        {Engenheira Mecânica Sênior}
        {AeroTech Engenharia de Sistemas}
        {\url{https://www.aerotech-engenharia.com.br}}
        {}
        {\begin{itemize}
\item Liderança técnica de uma equipe de 8 engenheiros e projetistas no desenvolvimento de sistemas térmicos e estruturais.
\item Responsável pela especificação, cálculo e validação de componentes críticos, atendendo requisitos de desempenho e segurança.
\item Implantação de fluxo de simulação CAE integrado ao CAD, reduzindo o ciclo de prototipagem em 30\%.
\item Gestão de fornecedores e transferência de projetos para a produção, com foco em qualidade e custo.
\item Aplicação de FMEA, análise de causa raiz e planos de ação para melhoria contínua.
\end{itemize}
\textbf{Palavras-chave}: Projeto mecânico, Simulação estrutural, ISO 9001, Manufatura enxuta, Gestão de equipes}

\cventry{abr. de 2018–dez. de 2020}
        {Engenheira de Projetos}
        {Metalúrgica Nova Era S.A.}
        {\url{https://www.metalurgicanovaera.com.br}}
        {}
        {\begin{itemize}
\item Desenvolvimento de produtos metálicos sob medida, desde o conceito até a liberação para fabricação.
\item Criação de modelos 3D, desenhos 2D e memoriais de cálculo com foco em manufaturabilidade.
\item Análise de falhas em campo e proposição de melhorias em projetos e processos.
\item Negociação com fornecedores e apoio à engenharia de processos na definição de tolerâncias.
\item Redução de 12\% no custo de matéria-prima por meio de revisão de projetos e padronização de componentes.
\end{itemize}
\textbf{Palavras-chave}: Desenvolvimento de produto, CAD 3D, Análise de falhas, Fornecedores, Redução de custos}

\cventry{jan. de 2016–mar. de 2018}
        {Engenheira Mecânica Júnior}
        {Hidráulica Vale Bombas e Sistemas}
        {\url{https://www.hidraulicavale.com.br}}
        {}
        {\begin{itemize}
\item Suporte ao projeto e à melhoria de bombas centrífugas e sistemas hidráulicos industriais.
\item Execução de ensaios de desempenho, análise de dados e elaboração de relatórios técnicos.
\item Atualização de desenhos e documentação técnica conforme normas internas e ABNT.
\item Apoio à manutenção e à equipe de campo na solução de problemas de instalação.
\end{itemize}
\textbf{Palavras-chave}: Bombas centrífugas, Ensaios de desempenho, Manutenção, Normas técnicas, CAD 2D}
\section{Idiomas}

\cvline{Português}{Nativo ou Bilíngue \hfill \textbf{Palavras-chave}: Nativo}
\cvline{Inglês}{Proficiência Profissional Fluente \hfill \textbf{Palavras-chave}: TOEFL iBT 105, Inglês técnico}
\cvline{Espanhol}{Proficiência Profissional Limitada \hfill \textbf{Palavras-chave}: Comunicação técnica}
\section{Competências}

\cvline{Projeto Mecânico e CAD}{Especialista \hfill \textbf{Palavras-chave}: SolidWorks, Siemens NX, CATIA V5, AutoCAD, GD\&T, Desenho técnico}
\cvline{Simulação e Análise Estrutural}{Avançado \hfill \textbf{Palavras-chave}: ANSYS Mechanical, Abaqus, FEA, CFD, MATLAB, Otimização topológica}
\cvline{Manufatura e Processos}{Avançado \hfill \textbf{Palavras-chave}: Usinagem CNC, Soldagem, Conformação de metais, Lean Manufacturing, Six Sigma, Metrologia}
\cvline{Gestão de Projetos e Qualidade}{Avançado \hfill \textbf{Palavras-chave}: PMBOK, FMEA, ISO 9001, APQP, SAP ERP, Minitab}
\cvline{Programação e Automação}{Intermediário \hfill \textbf{Palavras-chave}: Python, VBA, MATLAB/Simulink, Arduino, CLP Siemens, Git}
\section{Prêmios}

\cventry{nov. de 2022}
        {AeroTech Engenharia de Sistemas}
        {Prêmio de Inovação em Engenharia}
        {}
        {}
        {\begin{itemize}
\item Reconhecimento pelo desenvolvimento de um sistema de resfriamento modular com redução de 18\% no consumo energético.
\item Projeto conduzido em parceria com as áreas de manufatura, qualidade e suprimentos.
\end{itemize}}

\cventry{mar. de 2015}
        {Universidade Estadual de Campinas}
        {Menção Honrosa em Dissertação de Mestrado}
        {}
        {}
        {\begin{itemize}
\item Distinção pela qualidade técnica e aplicabilidade industrial da pesquisa em otimização topológica.
\end{itemize}}
\section{Certificados}

\cventry{ago. de 2019}
        {ASQ}
        {Certificação Lean Six Sigma Green Belt}
        {\url{https://asq.org/cert/lean-six-sigma-green-belt}}
        {}
        {}

\cventry{jun. de 2017}
        {Dassault Systèmes}
        {SolidWorks Professional}
        {\url{https://www.solidworks.com/certification}}
        {}
        {}

\cventry{set. de 2020}
        {PMI}
        {Curso de Gestão de Projetos - PMBOK}
        {\url{https://www.pmi.org/certifications/project-management-pmp}}
        {}
        {}
\section{Publicações}

\cventry{jul. de 2018}
        {Otimização Topológica de Componentes Estruturais para Veículos Comerciais}
        {Revista Brasileira de Engenharia Mecânica}
        {\url{https://www.rbem.org.br/artigo/otimizacao-topologica-veiculos}}
        {}
        {\begin{itemize}
\item Estudo comparativo entre métodos de otimização topológica aplicados a suportes estruturais.
\item Validação por elementos finitos e ensaios experimentais em protótipos metálicos.
\item Discussão sobre ganhos de massa, rigidez e custo de fabricação.
\end{itemize}}
\section{Projetos}

\cventry{fev. de 2022–nov. de 2023}
        {Projeto de trocador de calor compacto para reduzir consumo energético em linhas de produção.}
        {Sistema de Resfriamento Modular para Equipamentos Industriais}
        {\url{https://www.aerotech-engenharia.com.br/projetos/resfriamento-modular}}
        {}
        {\begin{itemize}
\item Coordenação do projeto desde a concepção até a validação em campo.
\item Simulação CFD para dimensionamento térmico e otimização do fluxo de ar.
\item Uso de manufatura aditiva na produção de componentes de teste.
\end{itemize}
\textbf{Palavras-chave}: Trocador de calor, CFD, Sustentabilidade, Manufatura aditiva}

\cventry{mai. de 2017–dez. de 2018}
        {Bancada automatizada para medição de vazão, pressão e eficiência de bombas.}
        {Plataforma de Ensaio de Bombas Centrífugas}
        {\url{https://www.hidraulicavale.com.br/projetos/plataforma-ensaio}}
        {}
        {\begin{itemize}
\item Especificação de instrumentação, aquisição de dados e painel de controle.
\item Padronização de ensaios conforme a norma ISO 9906.
\item Treinamento da equipe técnica para operação e manutenção da bancada.
\end{itemize}
\textbf{Palavras-chave}: Instrumentação, Aquisição de dados, Norma ISO 9906, Automação}
\section{Interesses}

\cvline{Esportes}{Corrida de rua, Ciclismo, Natação, Trilhas}
\cvline{Cultura e Tecnologia}{Leitura, Impressão 3D, Fotografia, Astronomia}
\section{Voluntariado}

\cventry{ago. de 2019–Atual}
        {Mentora de Projetos}
        {Instituto Engenheiros do Amanhã}
        {\url{https://www.engenheirosdoamanha.org.br}}
        {}
        {\begin{itemize}
\item Mentoria de estudantes de engenharia em projetos de mobilidade e eficiência energética.
\item Apoio na estruturação de cronogramas, orçamentos e apresentações técnicas.
\item Organização de workshops sobre CAD, simulação e boas práticas de projeto.
\end{itemize}}
    
\end{document}